% p109_1; source PDF page 109; manually verified editable table.
\begin{tikzpicture}[x=1cm,y=-1cm,inner sep=0pt,outer sep=0pt]
\path[use as bounding box] (0,0) rectangle (16.000,3.165);
\draw[line width=.4pt] (0.000,0.000) rectangle (5.200,0.855);
\node[align=center,text width=4.940cm,font={\linespread{1}\fontsize{10}{13.50}\selectfont\bfseries}] at (2.600,0.427) {责任能力的程度};
\draw[line width=.4pt] (5.200,0.000) rectangle (10.400,0.855);
\node[align=center,text width=4.940cm,font={\linespread{1}\fontsize{10}{13.50}\selectfont\bfseries}] at (7.800,0.427) {精神状态};
\draw[line width=.4pt] (10.400,0.000) rectangle (16.000,0.855);
\node[align=center,text width=5.340cm,font={\linespread{1}\fontsize{10}{13.50}\selectfont\bfseries}] at (13.200,0.427) {处理结论};
\draw[line width=.4pt] (0.000,0.855) rectangle (5.200,1.860);
\node[align=left,text width=4.940cm,font={\linespread{1}\fontsize{10}{13.50}\selectfont}] at (2.600,1.357) {完全责任能力};
\draw[line width=.4pt] (5.200,0.855) rectangle (10.400,1.860);
\node[align=left,text width=4.940cm,font={\linespread{1}\fontsize{10}{13.50}\selectfont}] at (7.800,1.357) {完全有辨认和控制能力};
\draw[line width=.4pt] (10.400,0.855) rectangle (16.000,1.860);
\node[align=left,text width=5.340cm,font={\linespread{1}\fontsize{10}{13.50}\selectfont}] at (13.200,1.357) {应负刑事责任};
\draw[line width=.4pt] (0.000,1.860) rectangle (5.200,3.165);
\node[align=left,text width=4.940cm,font={\linespread{1}\fontsize{10}{13.50}\selectfont}] at (2.600,2.513) {相对责任能力（又称限定或限制责任能力）};
\draw[line width=.4pt] (5.200,1.860) rectangle (10.400,3.165);
\node[align=left,text width=4.940cm,font={\linespread{1}\fontsize{10}{13.50}\selectfont}] at (7.800,2.513) {尚未完全丧失辨认和控制能力};
\draw[line width=.4pt] (10.400,1.860) rectangle (16.000,3.165);
\node[align=left,text width=5.340cm,font={\linespread{1}\fontsize{10}{13.50}\selectfont}] at (13.200,2.513) {应负刑事责任，可以从宽处罚};
\end{tikzpicture}
